% Generated draft: requires mathematical and source review.
\documentclass[11pt]{article}
\usepackage[margin=25mm]{geometry}
\usepackage{fontspec}
\setmainfont{Latin Modern Roman}
\usepackage{amsmath,amssymb,mathtools,unicode-math}
\setmathfont{Latin Modern Math}
\usepackage{enumitem,needspace,xurl,hyperref}
\setlength{\parindent}{0pt}
\setlength{\parskip}{5pt}
\setlength{\emergencystretch}{3em}
\newcommand{\N}{\mathbb{N}}
\newcommand{\Z}{\mathbb{Z}}
\newcommand{\Q}{\mathbb{Q}}
\newcommand{\R}{\mathbb{R}}
\newcommand{\C}{\mathbb{C}}
\newcommand{\F}{\mathbb{F}}
\newcommand{\A}{\mathbb{A}}
\newcommand{\PP}{\mathbb{P}}
\newcommand{\E}{\mathbb{E}}
\newcommand{\eps}{\varepsilon}
\newcommand{\dd}{\,\mathrm{d}}
\begin{document}
\section*{UnsolvedMath ID: 30003893}
\textbf{Source problem code:} catalogue entry
\subsection{Definitions and mathematical statement}
Let \(M\) be a closed smooth manifold, meaning a compact smooth manifold without boundary, and let \(g\) be a Riemannian metric satisfying
\[
\operatorname{scal}_g(x)>0\qquad\text{for every }x\in M.
\]
The source uses \(c\) to denote a delocalized continuous cyclic cocycle, \(\tau_c\) for an associated pairing functional, and \(\rho(M,g)\) for an associated secondary invariant. These symbols retain their source meaning; the supplied record does not specify the algebra and topology supporting \(c\), its degree or delocalization condition, the group containing \(\rho(M,g)\), or the construction and domain of \(\tau_c\).

\textbf{Requested computation.} Determine, in general, the pairing of delocalized continuous cyclic cocycles with the topological structure set. The source also describes this task as computing the associated expressions
\[
\tau_c\circ\rho(M,g).
\]
If \(\rho(M,g)\) denotes a class on which \(\tau_c\) is defined, the corresponding evaluation would be written
\[
\tau_c\bigl(\rho(M,g)\bigr).
\]
The request concerns the general admissible range of manifolds, metrics, and cocycles, rather than a specified example.

\textbf{Unresolved qualifications.} The supplied record does not define the topological structure set, its reference manifold or category, the relation between its elements and the positive scalar curvature data \((M,g)\), or the hypotheses ensuring that the secondary invariant and its pairing exist. In particular, no dimension, orientation, spin, fundamental-group, or degree-compatibility assumptions are stated. None are imposed here without evidence. The claimed equivalence between the structure-set pairing and the displayed secondary-invariant expressions is retained as a source assertion; the supplied data do not provide the constructions needed to formulate or justify that equivalence precisely. Consequently, this is a qualified statement of the requested computation, not a fully specified pairing problem.
\subsection{Short English statement}
For a closed manifold equipped with a positive scalar curvature metric, compute the general delocalized cyclic-cocycle pairing described by the source, including its asserted interpretation using the topological structure set. The required structure set, cocycle algebra, secondary invariant, and compatibility hypotheses are unspecified, so a precise unconditional formulation requires further source clarification.
\subsection{Sources}
\begin{itemize}
\item Supplied local UnsolvedMath registry record, permanent ID \(30003893\), title \emph{Delocalized Cyclic Pairings with Structure Sets}, source problem code \texttt{catalogue entry}, in the Oberwolfach Reports Open Problems collection. The record lists proposed year 2018 and supplies no individual proposer.
\item Supplied report citation: \emph{Non-commutative Geometry, Index Theory and Mathematical Physics} (2019), Oberwolfach Reports, pp.\ 1911--1981, DOI \texttt{10.4171/owr/2018/32}. \url{https://doi.org/10.4171/owr/2018/32}. The supplied citation uses 2019, while the proposed year and DOI identifier contain 2018; these source details are retained without reconciliation.
\item Supplied literature source: Xie--Yu Oberwolfach report. \url{https://ems.press/content/serial-article-files/46757?nt=1}.
\item Supplied literature source: Chen--Wang--Xie--Yu, \emph{Annals of K-Theory}. \url{https://doi.org/10.2140/akt.2023.8.531}.
\item Supplied literature source: Oberwolfach report. \url{https://doi.org/10.4171/owr/2018/32}.
\item Literature assessment dated exactly \texttt{2026-08-21}: the supplied curation labels the problem open and states that the general structure-set pairing is explicitly identified as an open problem, while later work computes important special cases. This is supplied dated curation derived from the OWR material, not an independently verified literature assessment. The record expressly calls for independent verification of time-sensitive status and cited claims.
\item The registry supplies the statement status \texttt{corrected\_verified} and reports manual semantic, typographic, and LaTeX review asserting self-containment and preservation of the mathematical target. That supplied review label does not resolve the formulation gaps recorded here or constitute independent mathematical verification.
\end{itemize}
\end{document}
